% Generated draft: requires mathematical and source review.
\documentclass[11pt]{article}
\usepackage[margin=25mm]{geometry}
\usepackage{fontspec}
\setmainfont{Latin Modern Roman}
\usepackage{amsmath,amssymb,mathtools,unicode-math}
\setmathfont{Latin Modern Math}
\usepackage{enumitem,needspace,xurl,hyperref}
\setlength{\parindent}{0pt}
\setlength{\parskip}{5pt}
\setlength{\emergencystretch}{3em}
\newcommand{\N}{\mathbb{N}}
\newcommand{\Z}{\mathbb{Z}}
\newcommand{\Q}{\mathbb{Q}}
\newcommand{\R}{\mathbb{R}}
\newcommand{\C}{\mathbb{C}}
\newcommand{\F}{\mathbb{F}}
\newcommand{\A}{\mathbb{A}}
\newcommand{\PP}{\mathbb{P}}
\newcommand{\E}{\mathbb{E}}
\newcommand{\eps}{\varepsilon}
\newcommand{\dd}{\,\mathrm{d}}
\begin{document}
\section*{UnsolvedMath ID: 30005145}
\textbf{Source problem code:} catalogue entry
\subsection{Definitions and mathematical statement}
An ancient mean-curvature flow of hypersurfaces in Euclidean space is a family \((M_t^n)_{t\in(-\infty,T)}\) in \(\R^{n+1}\), defined for all sufficiently negative times, whose smooth local parametrizations satisfy
\[
(\partial_t X)^\perp=-H\nu,
\]
where \(\nu\) is a chosen unit normal and \(H\) is the sum of the principal curvatures with the corresponding sign convention. The source's reference to flows in \(\R^4\) is understood as referring to three-dimensional hypersurfaces.

Under the usual mean-convex interpretation, an \(\alpha\)-noncollapsed flow has \(H>0\) and a constant \(\alpha>0\), independent of position and time, such that at every \(p\in M_t\) there are tangent interior and exterior balls, disjoint from \(M_t\) in their interiors and lying on the respective sides, each of radius at least
\[
\frac{\alpha}{H(p,t)}.
\]
Here the interior is the mean-convex side of a region bounded by the hypersurface. This is a conventional interpretation of ``noncollapsed''; the supplied excerpt does not specify the precise definition, the global region hypotheses, or whether degenerate cases with \(H=0\) are admitted.

For \(n\geq2\), write the principal curvatures as \(\lambda_1\leq\cdots\leq\lambda_n\). Strict two-convexity means \(\lambda_1+\lambda_2>0\); uniform two-convexity means that some fixed \(\beta>0\) satisfies \(\lambda_1+\lambda_2\geq\beta H\) everywhere in spacetime. The excerpt does not distinguish these conventions.

\textbf{Classification problem.} Classify ancient noncollapsed mean-curvature flows of hypersurfaces in higher-dimensional Euclidean spaces without imposing a two-convexity hypothesis. The absence of this hypothesis does not exclude flows that happen to be two-convex. The requested classification concerns ancient flows generally, without requiring self-similarity or the translating-soliton property.

\textbf{Unresolved scope.} The supplied record does not specify the exact range of dimensions, completeness, embeddedness, properness, connectedness, compactness, admissible singularities, or the equivalence relation for classification. Consequently, it does not determine a single fully specified classification theorem. The conventional definitions above explain the terminology and require confirmation against the complete mathematical source before they can fix the intended class.

The original excerpt describes classification in \(\R^4\) under a self-similarity assumption as a first step. Separately, the supplied dated curation reports a later classification of ancient noncollapsed flows in \(\R^4\) and labels the general higher-dimensional problem open. These are source-specific historical and curatorial claims, not independently verified conclusions of this draft.
\subsection{Short English statement}
Describe all ancient noncollapsed mean-curvature flows in higher dimensions without assuming two-convexity or self-similarity. The precise dimension range and global geometric hypotheses remain unspecified. The supplied curation labels the general problem open and reports later progress in \(\R^4\); these assessments require independent verification.
\subsection{Sources}
\begin{itemize}
\item Local UnsolvedMath registry: supplied record with permanent ID \(30005145\), exact source problem code \texttt{catalogue entry}, titled ``Classification of Ancient Noncollapsed Mean Curvature Flows,'' in the Oberwolfach Reports Open Problems collection. The registry labels the problem \texttt{open} and the statement \texttt{corrected\_verified}; these are supplied metadata, not mathematical verification by this draft.
\item Supplied primary-source citation: \emph{Geometrie (2023)}, Oberwolfach Reports, DOI 10.4171/owr/2022/28, pp.\ 1551--1601. \url{https://doi.org/10.4171/owr/2022/28}. The registry records proposed year 2022. Both year indications are retained as supplied.
\item The original excerpt contains the presentation heading ``Classification of noncollapsed translators in R4,'' attributed to Robert Haslhofer, with joint work credited to Kyeongsu Choi and Or Hershkovits. The registry does not name a proposer of the general classification problem.
\item Supplied literature source: Choi--Haslhofer, ``classification in R4.'' \url{https://arxiv.org/abs/2412.10581}.
\item Supplied literature source: Du--Haslhofer, ``recent higher-dimensional context.'' \url{https://msp.org/gt/2025/29-2/gt-v29-n2-p08-p.pdf}.
\item Supplied literature source: ``OWR report.'' \url{https://ems.press/content/serial-article-files/46965}.
\item Literature assessment dated exactly \texttt{2026-08-22}: the supplied curation states that classification without two-convexity remains unresolved in higher dimensions and that the 2024 work classifies ancient noncollapsed flows in \(\R^4\), while leaving the general higher-dimensional problem as a major challenge. This is supplied dated curation derived from the OWR material, not an independently verified literature assessment. Its time-sensitive status and cited claims require independent checking.
\end{itemize}
\end{document}
